\documentclass[11pt,a4paper]{amsart}

\usepackage[margin=1in]{geometry}
\usepackage{amsmath,amssymb,amsthm,mathtools}
\usepackage{microtype}
\usepackage{booktabs}
\usepackage[hidelinks]{hyperref}
\hypersetup{
  pdftitle={Integer Solutions of y2+x2y+xz2+2=0},
  pdfauthor={},
  pdfsubject={A complete classification of the integer solutions of a cubic Diophantine equation}
}

\numberwithin{equation}{section}
\allowdisplaybreaks

\newtheorem{theorem}{Theorem}[section]
\newtheorem{proposition}[theorem]{Proposition}
\newtheorem{lemma}[theorem]{Lemma}
\newtheorem{corollary}[theorem]{Corollary}
\newtheorem{remark}[theorem]{Remark}
\newtheorem{example}[theorem]{Example}

\theoremstyle{definition}
\newtheorem{definition}[theorem]{Definition}

\DeclareMathOperator{\sgn}{sgn}
\newcommand{\Z}{\mathbb Z}
\newcommand{\N}{\mathbb N}
\newcommand{\Q}{\mathbb Q}
\newcommand{\SL}{\operatorname{SL}}

\title[Integer solutions of a cubic Diophantine equation]
{Integer Solutions of \(y^2+x^2y+xz^2+2=0\)}
\author{Jamal Agbanwa}
\address{Independent Researcher}
\date{}

\begin{document}

\begin{abstract}
We determine all integer triples \((x,y,z)\) satisfying
\(y^2+x^2y+xz^2+2=0\).
The principal result is an exact divisor--square parametrization in which the parameter \(|x|\) is recovered uniquely from the solution.
We then give a complete fixed-\(|x|\) description: the positive-\(x\) branch is finite, the negative branch with square \(|x|\) is governed by a finite divisor calculation, and every solvable negative branch with nonsquare \(|x|\) is a finite disjoint union of explicit Pell recurrences.
We also prove the precise congruence obstruction on \(|x|\), a global unimodular parametrization, and an explicit infinite family with \(x=-2\).
\end{abstract}

\maketitle

\section{Statement of the classification}

We study the integral points on
\begin{equation}
 y^2+x^2y+xz^2+2=0,
 \qquad x,y,z\in\Z.
 \label{eq:E}
\end{equation}
Throughout, \(\N=\{1,2,3,\ldots\}\).  For \(m\in\N\) and
\(y\in\Z\) such that \(m\mid y^2+2\), define
\begin{equation}
 F_m(y):=my+\frac{y^2+2}{m}.
 \label{eq:Fdef}
\end{equation}

The following theorem is the basic complete parametrization.

\begin{theorem}[Exact divisor--square parametrization]\label{thm:main}
The integer solutions of \eqref{eq:E} are exactly the triples
\begin{equation}
 (x,y,z)=(-\varepsilon m,\,y,\,\nu t),
 \label{eq:mainparam}
\end{equation}
where
\begin{equation}
 m,t\in\N,\qquad y\in\Z,\qquad
 m\mid y^2+2,\qquad
 F_m(y)=\varepsilon t^2,
 \qquad \varepsilon,\nu\in\{\pm1\}.
 \label{eq:mainconditions}
\end{equation}
For every solution, the parameters \(m=|x|\), \(y\),
\(t=|z|\), and \(\varepsilon=-\sgn(x)\) are uniquely determined;
only \(\nu=\sgn(z)\) records the choice of sign of \(z\).
In particular, \(F_m(y)\) never vanishes for an admissible pair
\((m,y)\) occurring in a solution.
\end{theorem}

Thus an equivalent concise description is
\begin{equation}
 (x,y,z)=
 \left(-\sgn(F_m(y))m,\ y,\ \pm\sqrt{|F_m(y)|}\right),
 \label{eq:signparam}
\end{equation}
where \(m\mid y^2+2\) and \(|F_m(y)|\) is a perfect square.
The nonvanishing assertion in Theorem~\ref{thm:main} makes
\eqref{eq:signparam} unambiguous.

The later sections refine Theorem~\ref{thm:main} in three ways.
First, we characterize exactly the positive integers \(m\) for which
the necessary congruence \(y^2\equiv-2\pmod m\) is soluble.
Second, for each fixed \(m\) we convert the remaining square condition
into a positive-definite equation, a finite factorization, or a finite
union of Pell orbits according to the sign of \(x\) and whether \(m\)
is a square.  Third, we give a global parametrization by matrices in
\(\SL_2(\Z)\).

\section{Proof of the exact parametrization}

We begin by excluding the two degeneracies that could otherwise cause
ambiguity.

\begin{lemma}\label{lem:nonzero}
Every solution of \eqref{eq:E} satisfies \(x\ne0\) and \(z\ne0\).
\end{lemma}

\begin{proof}
If \(x=0\), then \eqref{eq:E} becomes \(y^2+2=0\), which has no
integer solution.

Suppose next that \(z=0\).  Since \(x\ne0\), equation \eqref{eq:E}
becomes
\begin{equation}
 y(y+x^2)=-2.
 \label{eq:z0factor}
\end{equation}
Because \(x^2>0\), the second factor is strictly larger than the first.
The only ordered pairs of integral factors of \(-2\) with this property
are \((-2,1)\) and \((-1,2)\).  In either case their difference is
\(3\), so \eqref{eq:z0factor} would force \(x^2=3\), impossible.
Therefore \(z\ne0\).
\end{proof}

\begin{proof}[Proof of Theorem~\ref{thm:main}]
Let \((x,y,z)\) satisfy \eqref{eq:E}.  By Lemma~\ref{lem:nonzero},
we may write
\begin{equation}
 x=\delta m,
 \qquad m=|x|\in\N,
 \qquad \delta\in\{\pm1\}.
 \label{eq:xdelta}
\end{equation}
Then \(x^2=m^2\), and \eqref{eq:E} is
\begin{equation}
 y^2+m^2y+\delta mz^2+2=0.
 \label{eq:mdeltaeq}
\end{equation}
Reduction modulo \(m\) gives
\begin{equation}
 y^2+2\equiv0\pmod m.
 \label{eq:necessarycong}
\end{equation}
Consequently
\begin{equation}
 q:=\frac{y^2+2}{m}
 \label{eq:qdef}
\end{equation}
is a positive integer.  Dividing \eqref{eq:mdeltaeq} by \(m\) yields
\begin{equation}
 q+my+\delta z^2=0.
 \label{eq:linearized}
\end{equation}
By \eqref{eq:Fdef}, this says
\begin{equation}
 F_m(y)=-\delta z^2.
 \label{eq:Ffromsolution}
\end{equation}
Set \(t=|z|\), \(\varepsilon=-\delta\), and
\(\nu=\sgn(z)\).  Lemma~\ref{lem:nonzero} gives \(t\in\N\), and
\eqref{eq:Ffromsolution} becomes \(F_m(y)=\varepsilon t^2\).
Moreover \(x=\delta m=-\varepsilon m\) and \(z=\nu t\), proving that
every solution occurs in \eqref{eq:mainparam}--\eqref{eq:mainconditions}.

Conversely, suppose the data in \eqref{eq:mainconditions} are given and
set \((x,y,z)\) by \eqref{eq:mainparam}.  Since \(x^2=m^2\) and
\(z^2=t^2\), we have
\begin{align*}
 y^2+x^2y+xz^2+2
 &=y^2+m^2y-\varepsilon mt^2+2\\
 &=m\left(\frac{y^2+2}{m}+my-\varepsilon t^2\right)\\
 &=m\bigl(F_m(y)-\varepsilon t^2\bigr)=0.
\end{align*}
Thus every parameter choice in the theorem gives a solution.

For a given solution, \(m=|x|\), the displayed value of \(y\), and
\(t=|z|\) are forced.  Equation \eqref{eq:Ffromsolution} then forces
\(\varepsilon=-\sgn(x)\), while \(\nu\) is exactly the sign of \(z\).
Finally, \(F_m(y)=\varepsilon t^2\ne0\) because \(t\in\N\).
\end{proof}

The two signs of \(x\) lead to useful specialized forms.

\begin{corollary}[The positive branch]\label{cor:positivebranch}
The solutions with \(x>0\) are exactly
\begin{equation}
 (x,y,z)=
 \left(m,-a,\,\pm\sqrt{ma-\frac{a^2+2}{m}}\right),
 \label{eq:positiveparam}
\end{equation}
where
\begin{equation}
 m\in\N,\qquad 1\le a\le m^2-1,
 \qquad m\mid a^2+2,
 \label{eq:positiveconditions1}
\end{equation}
and
\begin{equation}
 ma-\frac{a^2+2}{m}
 \label{eq:positiveSquare}
\end{equation}
is a positive perfect square.
For a fixed \(m\), the substitution
\begin{equation}
 a\longmapsto m^2-a
 \label{eq:ainvolution}
\end{equation}
preserves the value in \eqref{eq:positiveSquare}.
\end{corollary}

\begin{proof}
Put \(x=m>0\) in \eqref{eq:E}.  Then
\begin{equation}
 y(y+m^2)=-mz^2-2<0.
 \label{eq:yrangeproof}
\end{equation}
The two factors on the left differ by the positive integer \(m^2\), so
they have opposite signs exactly when
\(-m^2<y<0\).  Write \(a=-y\); then
\(1\le a\le m^2-1\).  Theorem~\ref{thm:main}, with
\(F_m(-a)<0\), gives
\begin{equation*}
 z^2=-F_m(-a)=ma-\frac{a^2+2}{m},
\end{equation*}
which proves the parametrization and the square condition.  Conversely,
these conditions are a specialization of Theorem~\ref{thm:main} and
therefore produce a solution.

Finally,
\begin{align*}
 m(m^2-a)-\frac{(m^2-a)^2+2}{m}
 &=m^3-ma-\left(m^3-2ma+\frac{a^2+2}{m}\right)\\
 &=ma-\frac{a^2+2}{m}.
\end{align*}
Also \((m^2-a)^2+2\equiv a^2+2\pmod m\), so divisibility is
preserved.  This proves the final assertion.
\end{proof}


\begin{example}[Some positive branches]\label{ex:positivebranches}
Corollary~\ref{cor:positivebranch} gives the following complete lists
for three small positive values of $x$:
\begin{align*}
 x=2&:\qquad (y,z)=(-2,\pm1),\\
 x=3&:\qquad (y,z)=(-2,\pm2),\ (-7,\pm2),\\
 x=9&:\qquad (y,z)=(-22,\pm12),\ (-59,\pm12).
\end{align*}
For instance, when $m=9$ the only admissible values of $a$ for which
\eqref{eq:positiveSquare} is a square are $a=22$ and $a=59$; the
involution $a\mapsto81-a$ interchanges them.
\end{example}

\begin{corollary}[The negative branch]\label{cor:negativebranch}
The solutions with \(x<0\) are exactly
\begin{equation}
 (x,y,z)=
 \left(-m,y,\,\pm\sqrt{my+\frac{y^2+2}{m}}\right),
 \label{eq:negativeparam}
\end{equation}
where \(m\in\N\), \(m\mid y^2+2\), and
\begin{equation}
 my+\frac{y^2+2}{m}
 \label{eq:negativeSquare}
\end{equation}
is a positive perfect square.
\end{corollary}

\begin{proof}
This is Theorem~\ref{thm:main} with \(\varepsilon=1\), equivalently
with \(x=-m\).
\end{proof}

The quadratic nature of \eqref{eq:E} in \(y\) gives an elementary
symmetry that is occasionally useful.

\begin{proposition}[Vieta involution]\label{prop:vieta}
If \((x,y,z)\) satisfies \eqref{eq:E}, then so does
\begin{equation}
 (x,-x^2-y,z).
 \label{eq:vieta}
\end{equation}
This transformation is an involution.  On the parameters of
Corollary~\ref{cor:positivebranch}, it is exactly
\(a\mapsto m^2-a\).
\end{proposition}

\begin{proof}
For fixed \(x,z\), equation \eqref{eq:E} is the monic quadratic
\begin{equation*}
 Y^2+x^2Y+(xz^2+2)=0.
\end{equation*}
The sum of its two roots is \(-x^2\).  Hence, if one root is \(y\), the
other is \(-x^2-y\).  Applying the transformation twice returns \(y\),
so it is an involution.  If \(x=m>0\) and \(y=-a\), the new value is
\(-m^2+a=-(m^2-a)\), as claimed.
\end{proof}

\section{The exact congruence obstruction on \texorpdfstring{$|x|$}{|x|}}

The divisibility condition in Theorem~\ref{thm:main} can itself be
classified completely.


For an odd prime $p$ and an integer $a$ not divisible by $p$, write
$\left(\frac ap\right)=1$ when $a$ is a quadratic residue modulo $p$
and $\left(\frac ap\right)=-1$ otherwise.

\begin{lemma}[Gauss's lemma]\label{lem:gauss}
Let $p$ be an odd prime, let $a\not\equiv0\pmod p$, and put
$n=(p-1)/2$.  For $1\le j\le n$, let $r_j$ be the unique integer with
\begin{equation*}
 -\frac p2<r_j<\frac p2,
 \qquad r_j\equiv aj\pmod p.
\end{equation*}
If exactly $N$ of the integers $r_j$ are negative, then
\begin{equation}
 \left(\frac ap\right)=(-1)^N.
 \label{eq:gausslemma}
\end{equation}
\end{lemma}

\begin{proof}
The absolute values $|r_1|,\ldots,|r_n|$ are distinct.  Indeed, if
$|r_i|=|r_j|$, then $ai\equiv\pm aj\pmod p$, hence
$i\equiv\pm j\pmod p$.  In the ranges $1\le i,j\le n$, the plus sign
forces $i=j$, while the minus sign would give $i+j=p$, impossible
because $i+j\le p-1$.  Therefore the $|r_j|$ are a permutation of
$1,2,\ldots,n$.  Taking products modulo $p$ gives
\begin{equation*}
 a^n n!
 \equiv\prod_{j=1}^n r_j
 =(-1)^N\prod_{j=1}^n|r_j|
 =(-1)^N n!\pmod p.
\end{equation*}
Since $p\nmid n!$, cancellation yields
$a^n\equiv(-1)^N\pmod p$.

It remains only to identify $a^n$ with the Legendre symbol.  Fermat's
little theorem follows because multiplication by $a$ permutes the
nonzero residue classes modulo $p$:
\begin{equation*}
 a^{p-1}(p-1)!\equiv(p-1)!\pmod p,
\end{equation*}
so $a^{p-1}\equiv1\pmod p$.  The $n$ nonzero squares modulo $p$ are
distinct and are roots of $X^n-1$; since a nonzero polynomial of
degree $n$ over a field has at most $n$ roots, these are all the roots.
Thus $a^n\equiv1\pmod p$ exactly when $a$ is a square.  If $a$ is not
a square, then $a^n\not\equiv1\pmod p$ while
$(a^n)^2=a^{p-1}\equiv1\pmod p$, so $a^n\equiv-1\pmod p$.
Consequently $a^n\equiv\left(\frac ap\right)\pmod p$.  Both sides of
\eqref{eq:gausslemma} belong to $\{\pm1\}$, and the lemma follows.
The same criterion also proves multiplicativity: for
$a,b\not\equiv0\pmod p$,
\begin{equation*}
 \left(\frac{ab}{p}\right)
 \equiv (ab)^n
 \equiv a^n b^n
 \equiv \left(\frac ap\right)\left(\frac bp\right)\pmod p,
\end{equation*}
and both sides are again in $\{\pm1\}$.
\end{proof}

\begin{lemma}[The residue of \(-2\) modulo an odd prime]\label{lem:minus2prime}
Let \(p\) be an odd prime.  Then
\begin{equation}
 u^2\equiv-2\pmod p
 \label{eq:minus2modp}
\end{equation}
has a solution if and only if
\begin{equation}
 p\equiv1\quad\text{or}\quad3\pmod8.
 \label{eq:presidueclasses}
\end{equation}
When it is soluble, it has exactly two roots modulo \(p\).
\end{lemma}

\begin{proof}
By Lemma~\ref{lem:gauss}, applied first to \(a=-1\) and then to
\(a=2\), we obtain
\begin{equation}
 \left(\frac{-1}{p}\right)=(-1)^{(p-1)/2},
 \qquad
 \left(\frac{2}{p}\right)=(-1)^{(p^2-1)/8}.
 \label{eq:supplementarylaws}
\end{equation}
For completeness, the second count is as follows.  The least positive
residues of \(2,4,\ldots,p-1\) exceed \(p/2\) precisely for the
integers \(j\) with \(p/4<j\le(p-1)/2\).  Their number is
\begin{equation*}
 \frac{p-1}{2}-\left\lfloor\frac p4\right\rfloor,
\end{equation*}
whose parity, checked in the four odd classes modulo \(8\), is the
parity of \((p^2-1)/8\).  The first formula is immediate because all
\((p-1)/2\) residues of \(-j\) are negative in least absolute form.

By multiplicativity of the Legendre symbol,
\begin{equation}
 \left(\frac{-2}{p}\right)
 =(-1)^{(p-1)/2+(p^2-1)/8}.
 \label{eq:minus2symbol}
\end{equation}
Checking the four odd residue classes modulo \(8\), the exponent in
\eqref{eq:minus2symbol} is even exactly for \(p\equiv1,3\pmod8\).
Thus \(-2\) is a quadratic residue exactly in those two classes.
If a root \(u\) exists, then \(-u\) is a distinct root because
\(p\ne2\), and a nonzero quadratic polynomial over the field
\(\Z/p\Z\) has at most two roots.  Hence there are exactly two.
\end{proof}

\begin{lemma}[Lifting roots through odd prime powers]\label{lem:lifting}
Let \(p\) be an odd prime.  Every solution of
\begin{equation}
 u^2\equiv-2\pmod{p^k},\qquad k\ge1,
 \label{eq:liftstart}
\end{equation}
lifts uniquely modulo \(p^{k+1}\) to a solution of the same congruence.
Consequently, if \(-2\) is a quadratic residue modulo \(p\), then
\(u^2\equiv-2\pmod{p^e}\) has exactly two roots for every \(e\ge1\).
\end{lemma}

\begin{proof}
Write \(u^2+2=p^kc\).  A lift has the form \(u+tp^k\), with
\(t\) considered modulo \(p\).  Since \(2k\ge k+1\),
\begin{align*}
 (u+tp^k)^2+2
 &\equiv p^k(c+2ut)\pmod{p^{k+1}}.
\end{align*}
The residue of \(u\) modulo \(p\) is nonzero, because
\(u^2\equiv-2\not\equiv0\pmod p\).  Hence \(2u\) is invertible modulo
\(p\), and there is a unique \(t\pmod p\) satisfying
\(c+2ut\equiv0\pmod p\).  This proves unique lifting.  Starting from
the two roots modulo \(p\) supplied by Lemma~\ref{lem:minus2prime}
and iterating gives exactly two roots modulo every \(p^e\).
\end{proof}

\begin{theorem}[Soluble moduli]\label{thm:moduli}
For \(m\in\N\), the congruence
\begin{equation}
 y^2\equiv-2\pmod m
 \label{eq:minus2modm}
\end{equation}
is soluble if and only if
\begin{equation}
 m=2^{\epsilon}\prod_{i=1}^r p_i^{e_i},
 \qquad
 \epsilon\in\{0,1\},
 \qquad
 p_i\equiv1\ \text{or}\ 3\pmod8,
 \label{eq:mclassification}
\end{equation}
where the \(p_i\) are distinct odd primes and \(e_i\in\N\).
For such an \(m\), the congruence has exactly \(2^r\) residue classes
of solutions modulo \(m\).
\end{theorem}

\begin{proof}
Suppose first that \eqref{eq:minus2modm} is soluble.  If \(4\mid m\),
then it would imply \(y^2\equiv2\pmod4\), impossible because a square
modulo \(4\) is \(0\) or \(1\).  Thus the exponent of \(2\) in \(m\)
is at most one.  If an odd prime \(p\mid m\), reduction modulo \(p\)
shows that \(-2\) is a quadratic residue modulo \(p\).  By
Lemma~\ref{lem:minus2prime}, \(p\equiv1\) or \(3\pmod8\).  This proves
the necessity of \eqref{eq:mclassification}.

Conversely, assume \eqref{eq:mclassification}.  For each odd prime
power \(p_i^{e_i}\), Lemmas~\ref{lem:minus2prime} and
\ref{lem:lifting} provide exactly two roots.  Modulo \(2\), the unique
root is \(0\).  The moduli \(2^{\epsilon},p_1^{e_1},\ldots,p_r^{e_r}\)
are pairwise coprime, so the Chinese remainder theorem combines an
arbitrary choice of one root for each odd prime power into a unique
root modulo \(m\).  There are therefore exactly \(2^r\) roots.

For completeness, the Chinese remainder theorem in this setting
follows directly from Bezout's identity: for pairwise coprime moduli
\(n_j\), put \(N=\prod_j n_j\), \(N_j=N/n_j\), and choose
\(c_jN_j\equiv1\pmod{n_j}\); then
\(\sum_j a_jc_jN_j\) has the prescribed residue \(a_j\) modulo each
\(n_j\), and uniqueness holds modulo \(N\).
\end{proof}

\begin{corollary}[Restrictions on \(x,y,z\)]\label{cor:parity}
If \((x,y,z)\) satisfies \eqref{eq:E}, then
\begin{equation}
 |x|=2^{\epsilon}\prod_{i=1}^r p_i^{e_i}
 \label{eq:xrestriction}
\end{equation}
with the restrictions in Theorem~\ref{thm:moduli}.  Moreover,
\begin{equation}
 \begin{aligned}
 |x|\text{ odd} &\quad\Longrightarrow\quad z\text{ even},\\
 |x|\text{ even}&\quad\Longrightarrow\quad y\text{ even and }z\text{ odd}.
 \end{aligned}
 \label{eq:parityconclusions}
\end{equation}
\end{corollary}

\begin{proof}
The factorization restriction follows from
\eqref{eq:necessarycong} and Theorem~\ref{thm:moduli}.

If \(|x|\) is odd, then reduction of \eqref{eq:E} modulo \(2\) gives
\begin{equation*}
 y^2+y+z^2\equiv0\pmod2.
\end{equation*}
Since \(y^2+y\) is always even, \(z^2\) is even, and hence \(z\) is
even.

If \(|x|\) is even, Theorem~\ref{thm:moduli} shows that its exact
\(2\)-adic valuation is one, so \(x\equiv2\pmod4\).  Congruence
\eqref{eq:necessarycong} modulo \(2\) forces \(y\) to be even.
Reducing \eqref{eq:E} modulo \(4\), the terms \(y^2\) and \(x^2y\)
vanish, and we obtain
\begin{equation*}
 2z^2+2\equiv0\pmod4.
\end{equation*}
Thus \(z^2\equiv1\pmod2\), so \(z\) is odd.
\end{proof}

\section{Reduction to quadratic norm equations}

The remaining square condition in Theorem~\ref{thm:main} has a useful
symmetric form.

\begin{proposition}[Quadratic norm form]\label{prop:normform}
Let \(m\in\N\), \(\delta\in\{\pm1\}\), and put \(x=\delta m\).
There is a bijection between solutions \((x,y,z)\) of \eqref{eq:E}
with this fixed value of \(x\) and integer pairs \((U,V)\) satisfying
\begin{equation}
 U^2+\delta mV^2=m^4-8,
 \label{eq:normequation}
\end{equation}
together with the parity conditions
\begin{equation}
 U\equiv m^2\pmod2,
 \qquad V\equiv0\pmod2.
 \label{eq:normparity}
\end{equation}
The bijection is
\begin{equation}
 U=2y+m^2,
 \qquad V=2z,
 \label{eq:UVforward}
\end{equation}
and
\begin{equation}
 y=\frac{U-m^2}{2},
 \qquad z=\frac V2.
 \label{eq:UVbackward}
\end{equation}
\end{proposition}

\begin{proof}
With \(x=\delta m\), equation \eqref{eq:E} is
\begin{equation*}
 y^2+m^2y+\delta mz^2+2=0.
\end{equation*}
Multiplying by \(4\) and completing the square in \(y\) gives
\begin{equation*}
 (2y+m^2)^2+\delta m(2z)^2=m^4-8,
\end{equation*}
which is \eqref{eq:normequation}.  The definitions in
\eqref{eq:UVforward} imply \eqref{eq:normparity}.

Conversely, the parity conditions make the quantities in
\eqref{eq:UVbackward} integers.  Substituting them into
\eqref{eq:normequation} and reversing the completion of the square
gives \eqref{eq:E}.  The two constructions are inverse to one another.
\end{proof}

\subsection{Positive \texorpdfstring{\(x\)}{x}}

For \(x=m>0\), the norm equation is positive definite.

\begin{corollary}[Finiteness for fixed positive \(x\)]\label{cor:positivefinite}
For every fixed positive integer \(m\), there are only finitely many
solutions with \(x=m\).  More precisely, there are no such solutions
for \(m=1\), and for \(m\ge2\) every corresponding pair \((U,V)\)
satisfies
\begin{equation}
 |U|\le\sqrt{m^4-8},
 \qquad
 |V|\le\sqrt{\frac{m^4-8}{m}}.
 \label{eq:positivebounds}
\end{equation}
Checking the finitely many pairs in this box subject to
\eqref{eq:normparity} gives all solutions.
\end{corollary}

\begin{proof}
When \(m=1\), equation \eqref{eq:normequation} is
\(U^2+V^2=-7\), impossible.  For \(m\ge2\), it is
\begin{equation*}
 U^2+mV^2=m^4-8>0.
\end{equation*}
Both summands are nonnegative, giving the two bounds in
\eqref{eq:positivebounds}.  A bounded subset of \(\Z^2\) is finite,
and Proposition~\ref{prop:normform} completes the proof.
\end{proof}

\subsection{Negative \texorpdfstring{\(x\)}{x} with square \texorpdfstring{\(|x|\)}{|x|}}

When \(x=-m\) and \(m\) is a square, the norm equation factors over
\(\Z\).

\begin{theorem}[Finite divisor parametrization for \(x=-d^2\)]
\label{thm:squarem}
Fix \(d\in\N\).  The solutions of \eqref{eq:E} with \(x=-d^2\) are
in bijection with ordered pairs \((A,B)\in\Z^2\) satisfying
\begin{equation}
 AB=d^8-8,
 \label{eq:ABproduct}
\end{equation}
\begin{equation}
 A+B\equiv2d^4\pmod4,
 \qquad
 B-A\equiv0\pmod{4d}.
 \label{eq:ABcongruences}
\end{equation}
The corresponding solution is
\begin{equation}
 (x,y,z)=
 \left(-d^2,\ \frac{A+B-2d^4}{4},\ \frac{B-A}{4d}\right).
 \label{eq:ABsolution}
\end{equation}
In particular, for each square value of \(|x|\), there are only
finitely many solutions.
\end{theorem}

\begin{proof}
Set \(m=d^2\) and \(\delta=-1\) in
Proposition~\ref{prop:normform}.  The norm equation becomes
\begin{equation}
 U^2-d^2V^2=d^8-8.
 \label{eq:squarefactorstart}
\end{equation}
Define
\begin{equation}
 A=U-dV,
 \qquad B=U+dV.
 \label{eq:ABdef}
\end{equation}
Then \eqref{eq:squarefactorstart} is exactly
\eqref{eq:ABproduct}.  Since \(U=2y+d^4\) and \(V=2z\), we also have
\begin{equation*}
 A+B=4y+2d^4,
 \qquad B-A=4dz,
\end{equation*}
which proves the congruences in \eqref{eq:ABcongruences} and the
recovery formulas in \eqref{eq:ABsolution}.

Conversely, suppose \((A,B)\) satisfies
\eqref{eq:ABproduct}--\eqref{eq:ABcongruences}.  Put
\begin{equation*}
 U=\frac{A+B}{2},
 \qquad V=\frac{B-A}{2d}.
\end{equation*}
The first congruence in \eqref{eq:ABcongruences} implies that \(U\) is
an integer with \(U\equiv d^4=m^2\pmod2\); the second implies that
\(V\) is an even integer.  Equation \eqref{eq:ABproduct} gives
\(U^2-d^2V^2=d^8-8\).  Proposition~\ref{prop:normform} therefore
produces the solution \eqref{eq:ABsolution}.  The two constructions are
inverse.  Finally, the nonzero integer \(d^8-8\) has finitely many
ordered factor pairs, proving finiteness.
\end{proof}

\begin{example}
For \(d=1\), Theorem~\ref{thm:squarem} gives exactly
\begin{equation}
 (x,y,z)=(-1,1,\pm2),
 \qquad
 (-1,-2,\pm2).
 \label{eq:xminusone}
\end{equation}
Indeed, in this case \((A,B)\) runs through the admissible ordered
factor pairs of \(-7\).
\end{example}

\subsection{Negative \texorpdfstring{\(x\)}{x} with nonsquare \texorpdfstring{\(|x|\)}{|x|}}

We first prove the elementary Pell fact needed to make the description
self-contained.

\begin{lemma}[Existence of a positive Pell unit]\label{lem:pell}
If \(m\in\N\) is not a square, then there exist \(a,b\in\N\) such that
\begin{equation}
 a^2-mb^2=1.
 \label{eq:pellunit}
\end{equation}
\end{lemma}

\begin{proof}
For every integer \(H\ge2\), apply the pigeonhole principle to the
\(H+1\) fractional parts
\begin{equation*}
 0,\{\sqrt m\},\{2\sqrt m\},\ldots,\{H\sqrt m\}
\end{equation*}
placed in \(H\) half-open intervals of length \(1/H\).  Two lie in the
same interval.  Subtracting their indices gives integers \(p,q\) with
\(1\le q\le H\) and
\begin{equation}
 |q\sqrt m-p|<\frac1H.
 \label{eq:dirichlet}
\end{equation}
Because \(m\) is not a square,
\begin{equation}
 c:=p^2-mq^2
 \label{eq:cdef}
\end{equation}
is a nonzero integer.  Also
\begin{align}
 |c|
 &=|p-q\sqrt m|\,|p+q\sqrt m|\notag\\
 &<\frac1H\left(2H\sqrt m+1\right)
 <2\sqrt m+1.
 \label{eq:cbound}
\end{align}
The pairs \((p,q)\) obtained as \(H\) varies cannot belong to a finite
set: if they did, the positive numbers \(|q\sqrt m-p|\) arising from
that finite set would have a positive minimum, contradicting
\eqref{eq:dirichlet} for sufficiently large \(H\).  Thus infinitely
many distinct pairs occur.  Since \eqref{eq:cbound} leaves only
finitely many possible nonzero values of \(c\), one fixed value of
\(c\) occurs for infinitely many pairs.  Passing to a further infinite
subsequence, we may assume that \(p\) and \(q\) have fixed residue
classes modulo \(|c|\).

For all sufficiently large members of this subsequence, \(p>0\).
Choose two distinct pairs \((p_1,q_1)\), \((p_2,q_2)\) from it such
that
\begin{equation*}
 p_2+q_2\sqrt m>p_1+q_1\sqrt m>0.
\end{equation*}
Because the two pairs are congruent componentwise modulo \(|c|\), the
integers
\begin{equation}
 A=\frac{p_2p_1-mq_2q_1}{c},
 \qquad
 B=\frac{q_2p_1-p_2q_1}{c}
 \label{eq:ABpell}
\end{equation}
are well defined.  Indeed, modulo \(c\) the first numerator is
\(p_1^2-mq_1^2=c\), and the second is zero.  In the real quadratic
field,
\begin{equation}
 A+B\sqrt m
 =\frac{p_2+q_2\sqrt m}{p_1+q_1\sqrt m}>1.
 \label{eq:pellratio}
\end{equation}
The norm of this quotient is \(c/c=1\).  Its conjugate is therefore
the positive number \((A+B\sqrt m)^{-1}<1\).  Hence
\begin{equation*}
 A=\frac{(A+B\sqrt m)+(A-B\sqrt m)}2>0,
\end{equation*}
and
\begin{equation*}
 B\sqrt m
 =\frac{(A+B\sqrt m)-(A-B\sqrt m)}2>0.
\end{equation*}
Thus \(A,B\in\N\), and taking norms in \eqref{eq:pellratio} gives
\(A^2-mB^2=1\).  Renaming \((A,B)\) as \((a,b)\) proves the lemma.
\end{proof}

\begin{theorem}[Complete Pell-orbit description]\label{thm:pellorbits}
Fix a nonsquare integer \(m\ge2\), and set
\begin{equation}
 N_m=m^4-8,
 \qquad R_m=\sqrt{N_m}.
 \label{eq:NmRm}
\end{equation}
Choose \(a,b\in\N\) with \(a^2-mb^2=1\), and write
\begin{equation}
 \eta=(a+b\sqrt m)^2=P+Q\sqrt m.
 \label{eq:etadef}
\end{equation}
Then
\begin{equation}
 P^2-mQ^2=1,
 \qquad P\text{ is odd},
 \qquad Q\text{ is even},
 \qquad \eta>1.
 \label{eq:etaProperties}
\end{equation}
Define the seed set
\begin{align}
 \mathcal B_m:=\bigl\{(U,V)\in\Z^2:\;&
 U^2-mV^2=N_m,\quad
 U\equiv m^2\pmod2,\quad V\equiv0\pmod2,\notag\\
 &R_m\le |U+V\sqrt m|<\eta R_m\bigr\}.
 \label{eq:seedset}
\end{align}
Then \(\mathcal B_m\) is finite, and every solution of \eqref{eq:E}
with \(x=-m\) is obtained uniquely by choosing a seed
\((U_0,V_0)\in\mathcal B_m\), an integer \(n\in\Z\), and setting
\begin{equation}
 U_n+V_n\sqrt m=(U_0+V_0\sqrt m)\eta^n,
 \label{eq:pellorbit}
\end{equation}
followed by
\begin{equation}
 (x,y,z)=\left(-m,\frac{U_n-m^2}{2},\frac{V_n}{2}\right).
 \label{eq:pellback}
\end{equation}
Equivalently, consecutive elements of each orbit satisfy
\begin{equation}
 \begin{aligned}
 U_{n+1}&=PU_n+mQV_n,\\
 V_{n+1}&=QU_n+PV_n.
 \end{aligned}
 \label{eq:pellrecurrence}
\end{equation}
Consequently, for every fixed nonsquare \(m\), the equation with
\(x=-m\) has either no integer solutions or infinitely many integer
solutions.
\end{theorem}

\begin{proof}
Lemma~\ref{lem:pell} supplies \(a,b\).  Expanding
\eqref{eq:etadef} gives
\begin{equation*}
 P=a^2+mb^2=1+2mb^2,
 \qquad Q=2ab.
\end{equation*}
Thus \(P\) is odd, \(Q\) is even, and \(P,Q>0\).  The norm of
\(\eta\) is \((a^2-mb^2)^2=1\), proving
\eqref{eq:etaProperties}.

We next prove that \(\mathcal B_m\) is finite.  For a seed put
\begin{equation*}
 \alpha=U+V\sqrt m,
 \qquad \overline\alpha=U-V\sqrt m.
\end{equation*}
Its norm is \(N_m>0\).  Since \(|\alpha|\ge R_m\),
\begin{equation*}
 |\overline\alpha|=\frac{N_m}{|\alpha|}\le R_m.
\end{equation*}
Together with \(|\alpha|<\eta R_m\), this gives
\begin{equation}
 |U|<\frac{\eta+1}{2}R_m,
 \qquad
 |V|<\frac{\eta+1}{2\sqrt m}R_m.
 \label{eq:seedbounds}
\end{equation}
Indeed, \(2U=\alpha+\overline\alpha\) and
\(2V\sqrt m=\alpha-\overline\alpha\).  The bounds
\eqref{eq:seedbounds} place all seeds in a bounded rectangle in
\(\Z^2\), so there are finitely many.

Multiplication by \(\eta\) preserves the norm equation because
\(P^2-mQ^2=1\).  In coordinates it is exactly the recurrence
\eqref{eq:pellrecurrence}.  Since \(P\) is odd and \(Q\) is even,
\begin{equation*}
 U_{n+1}\equiv U_n\pmod2,
 \qquad V_{n+1}\equiv V_n\pmod2.
\end{equation*}
The inverse unit is \(\eta^{-1}=P-Q\sqrt m\), so the same conclusions
hold for negative powers.  Thus every orbit in \eqref{eq:pellorbit}
consists entirely of parity-compatible solutions of the norm equation.
Proposition~\ref{prop:normform} then gives \eqref{eq:pellback}.

Conversely, let \((U,V)\) be any parity-compatible solution of
\(U^2-mV^2=N_m\), and put \(\alpha=U+V\sqrt m\ne0\).  Since
\(\eta>1\), there is a unique integer \(k\) such that
\begin{equation}
 R_m\le |\alpha|\eta^k<\eta R_m.
 \label{eq:fundamentalannulus}
\end{equation}
For example, one may take
\begin{equation*}
 k=-\left\lfloor\log_{\eta}(|\alpha|/R_m)\right\rfloor,
\end{equation*}
for then
\(0\le\log_{\eta}(|\alpha|/R_m)+k<1\), which is exactly
\eqref{eq:fundamentalannulus}.
The element \(\alpha\eta^k\) still belongs to \(\Z[\sqrt m]\), has
norm \(N_m\), and satisfies the required parity conditions.  It
therefore corresponds to a unique seed in \(\mathcal B_m\).  Reversing
the power shows that \((U,V)\) lies in the orbit of that seed.  The
half-open interval in \eqref{eq:fundamentalannulus} also proves
uniqueness of the seed and exponent.

If \(\mathcal B_m\) is empty, there are no solutions.  If it contains
a seed \(\alpha_0\), the elements \(\alpha_0\eta^n\), \(n\in\Z\),
are pairwise distinct because \(\eta>1\) and \(\alpha_0\ne0\).
Hence there are infinitely many solutions.
\end{proof}

\section{A global unimodular parametrization}

The divisor--square parametrization can be encoded by the unique
positive-definite integral binary quadratic form of discriminant
\(-8\).  We include the reduction argument in full.

\begin{lemma}[Forms of discriminant \(-8\)]\label{lem:formreduction}
Every positive-definite integral binary quadratic form
\begin{equation}
 f(S,T)=AS^2+BST+CT^2
 \label{eq:generalform}
\end{equation}
with discriminant
\begin{equation}
 B^2-4AC=-8
 \label{eq:discminus8}
\end{equation}
is properly equivalent over \(\Z\) to
\begin{equation}
 S^2+2T^2.
 \label{eq:principalform}
\end{equation}
That is, some matrix in \(\SL_2(\Z)\) transforms one form into the
other.
\end{lemma}

\begin{proof}
Because \(f\) is positive definite, it assumes a least positive value
on \(\Z^2\setminus\{(0,0)\}\); denote this value by \(A_0\), attained
at a vector \(e_1\).  The vector \(e_1\) is primitive, since otherwise
\(e_1=de_1'\) with \(|d|>1\) would give
\(f(e_1')=f(e_1)/d^2<A_0\).  By Bezout's identity, \(e_1\) extends to
an oriented integral basis \((e_1,e_2)\) of \(\Z^2\).

In this basis the form is
\begin{equation*}
 A_0S^2+B_0ST+C_0T^2.
\end{equation*}
Replace \(e_2\) by \(e_2+ke_1\), choosing \(k\in\Z\) nearest to
\(-B_0/(2A_0)\).  This basis change has determinant one, leaves the
first coefficient equal to \(A_0\), and changes the middle coefficient
to
\begin{equation*}
 B_1=B_0+2kA_0
\end{equation*}
with \(|B_1|\le A_0\).  Let the new last coefficient be \(C_1\).
Since the new second basis vector is nonzero and integral, the
minimality of \(A_0\) gives \(C_1\ge A_0\).  The discriminant remains
\(-8\), so
\begin{equation*}
 8=4A_0C_1-B_1^2
 \ge4A_0^2-A_0^2=3A_0^2.
\end{equation*}
As \(A_0\) is a positive integer, this forces \(A_0=1\).  Equation
\eqref{eq:discminus8} also implies that \(B_1\) is even.  Since
\(|B_1|\le1\), we get \(B_1=0\), and then
\(4C_1=8\), so \(C_1=2\).  All basis changes used had determinant
one, proving proper equivalence to \eqref{eq:principalform}.
\end{proof}

\begin{theorem}[Unimodular parametrization]\label{thm:SL2}
Let
\begin{equation}
 G=\begin{pmatrix}r&u\\ s&v\end{pmatrix}\in\SL_2(\Z),
 \qquad rv-su=1,
 \label{eq:Gdef}
\end{equation}
and define
\begin{equation}
 M=r^2+2s^2,
 \qquad Y=ru+2sv,
 \qquad Q=u^2+2v^2,
 \label{eq:MYQ}
\end{equation}
and
\begin{equation}
 T=MY+Q.
 \label{eq:Tdef}
\end{equation}
Then
\begin{equation}
 MQ-Y^2=2
 \label{eq:detidentity}
\end{equation}
and \(T\ne0\).  Whenever \(|T|=t^2\) is a perfect square, necessarily
\(t\in\N\), and
\begin{equation}
 (x,y,z)=\bigl(-\sgn(T)M,\,Y,\,\pm t\bigr)
 \label{eq:SL2solution}
\end{equation}
is a solution of \eqref{eq:E}.

Conversely, every integer solution of \eqref{eq:E} is obtained from at
least one matrix \(G\in\SL_2(\Z)\) in this way.
\end{theorem}

\begin{proof}
Expanding the left side of \eqref{eq:detidentity} gives
\begin{align*}
 MQ-Y^2
 &=(r^2+2s^2)(u^2+2v^2)-(ru+2sv)^2\\
 &=2(r^2v^2+s^2u^2-2rsuv)\\
 &=2(rv-su)^2=2.
\end{align*}
Suppose for contradiction that \(T=MY+Q=0\).  Since \(M,Q>0\), we
have \(Y<0\).  Substitution of \(Q=-MY\) into
\eqref{eq:detidentity} gives
\begin{equation*}
 -Y(M^2+Y)=2.
\end{equation*}
Writing \(a=-Y>0\), this becomes
\begin{equation*}
 a(M^2-a)=2.
\end{equation*}
Both factors are positive integers.  They must be \((1,2)\) or
\((2,1)\), and either case gives \(M^2=3\), impossible.  Hence
\(T\ne0\).

Now suppose \(|T|=t^2\).  Since \(T\ne0\), \(t\in\N\).  For either
choice of the sign of \(z\), set \((x,y,z)\) as in
\eqref{eq:SL2solution}.  Then \(x^2=M^2\), \(z^2=t^2\), and
\(\sgn(T)t^2=T\).  Therefore
\begin{align*}
 y^2+x^2y+xz^2+2
 &=Y^2+M^2Y-MT+2\\
 &=Y^2+M^2Y-M(MY+Q)+2\\
 &=Y^2-MQ+2=0
\end{align*}
by \eqref{eq:detidentity}.

For completeness, let \((x,y,z)\) be any solution.  Put
\begin{equation*}
 m=|x|,
 \qquad q=\frac{y^2+2}{m}.
\end{equation*}
Theorem~\ref{thm:main} shows that \(q\in\N\), and
\begin{equation}
 mq-y^2=2.
 \label{eq:mqidentity}
\end{equation}
Consider the positive-definite integral form
\begin{equation}
 f(S,T)=mS^2+2yST+qT^2.
 \label{eq:solutionform}
\end{equation}
Its discriminant is \(4y^2-4mq=-8\).  By
Lemma~\ref{lem:formreduction}, there is a matrix
\(G\in\SL_2(\Z)\) such that
\begin{equation}
 f(S,T)=(rS+uT)^2+2(sS+vT)^2.
 \label{eq:formmatrix}
\end{equation}
Comparing coefficients in \eqref{eq:formmatrix} gives
\begin{equation*}
 m=M,
 \qquad y=Y,
 \qquad q=Q.
\end{equation*}
The linearized equation \eqref{eq:linearized} says
\begin{equation*}
 T=MY+Q=my+q=-\sgn(x)z^2.
\end{equation*}
Thus \(|T|=z^2\) and
\(-\sgn(T)M=x\).  The given solution is therefore recovered from
\eqref{eq:SL2solution}.
\end{proof}

The dependence on the second column of \(G\) can be written as one
explicit quadratic polynomial.

\begin{corollary}[One-variable form of the square condition]
\label{cor:Tk}
Fix a primitive pair \((r,s)\in\Z^2\), choose one pair
\((u_0,v_0)\in\Z^2\) with
\begin{equation}
 rv_0-su_0=1,
 \label{eq:bezoutcolumn}
\end{equation}
and put
\begin{equation}
 M=r^2+2s^2,
 \quad A=ru_0+2sv_0,
 \quad Q_0=u_0^2+2v_0^2.
 \label{eq:MAQ0}
\end{equation}
Every second column \((u,v)^T\) completing \((r,s)^T\) to a matrix in
\(\SL_2(\Z)\) is uniquely of the form
\begin{equation}
 u=u_0+kr,
 \qquad v=v_0+ks,
 \qquad k\in\Z.
 \label{eq:secondcolumns}
\end{equation}
For this matrix, the signed square \(T\) in
Theorem~\ref{thm:SL2} is
\begin{equation}
 T(k)=Mk^2+(M^2+2A)k+MA+Q_0,
 \label{eq:Tk}
\end{equation}
and the discriminant of this quadratic polynomial is
\begin{equation}
 \operatorname{disc}T=M^4-8.
 \label{eq:Tkdisc}
\end{equation}
\end{corollary}

\begin{proof}
If another pair \((u,v)\) satisfies \(rv-su=1\), subtracting
\eqref{eq:bezoutcolumn} gives
\begin{equation*}
 r(v-v_0)=s(u-u_0).
\end{equation*}
Since \(\gcd(r,s)=1\), there is a unique \(k\in\Z\) with
\(u-u_0=kr\) and \(v-v_0=ks\).  This proves
\eqref{eq:secondcolumns}.

For these values,
\begin{equation*}
 Y=A+kM,
 \qquad Q=Q_0+2kA+k^2M.
\end{equation*}
Hence
\begin{align*}
 T(k)=MY+Q
 &=M(A+kM)+Q_0+2kA+k^2M\\
 &=Mk^2+(M^2+2A)k+MA+Q_0,
\end{align*}
which is \eqref{eq:Tk}.  Finally, the identity
\(MQ_0-A^2=2\), which is \eqref{eq:detidentity} for the initial
matrix, gives
\begin{align*}
 \operatorname{disc}T
 &=(M^2+2A)^2-4M(MA+Q_0)\\
 &=M^4+4(A^2-MQ_0)=M^4-8.
\end{align*}
\end{proof}

\section{An explicit infinite family}

The preceding classification shows structurally when a fixed negative
value of \(x\) can support infinitely many solutions.  The smallest
particularly transparent example is \(x=-2\).

\begin{theorem}[An infinite family with \(x=-2\)]\label{thm:infinitefamily}
Define positive integers \(A_n,B_n\) by
\begin{equation}
 A_n+B_n\sqrt2=(2+\sqrt2)(3+2\sqrt2)^n,
 \qquad n\ge0.
 \label{eq:AnBn}
\end{equation}
Then for every \(n\ge0\), all four triples
\begin{equation}
 (-2,A_n-2,\pm B_n),
 \qquad
 (-2,-A_n-2,\pm B_n)
 \label{eq:explicitfamily}
\end{equation}
solve \eqref{eq:E}.  These give infinitely many distinct integer
solutions.
\end{theorem}

\begin{proof}
Both factors on the right side of \eqref{eq:AnBn} belong to
\(\Z[\sqrt2]\), so \(A_n,B_n\in\Z\); they are positive because all
coefficients in the product are positive.  Taking norms gives
\begin{equation}
 A_n^2-2B_n^2
 =\bigl(2^2-2\cdot1^2\bigr)
  \bigl(3^2-2\cdot2^2\bigr)^n
 =2.
 \label{eq:AnBnnorm}
\end{equation}
For \(x=-2\), equation \eqref{eq:E} is equivalent to
\begin{equation}
 (y+2)^2-2z^2=2.
 \label{eq:xminus2pell}
\end{equation}
Taking \(y+2=\pm A_n\) and \(z=\pm B_n\) in
\eqref{eq:AnBnnorm} proves \eqref{eq:explicitfamily}.

Moreover, multiplication by \(3+2\sqrt2\) yields
\begin{equation}
 A_{n+1}=3A_n+4B_n,
 \qquad B_{n+1}=2A_n+3B_n.
 \label{eq:AnBnrec}
\end{equation}
In particular, \(A_{n+1}>A_n\), so the displayed solutions are
distinct as \(n\) varies.  Hence there are infinitely many.
\end{proof}

The first few values are
\begin{equation*}
 (A_n,B_n)=(2,1),(10,7),(58,41),(338,239),\ldots,
\end{equation*}
and hence the family begins
\begin{align*}
 &(-2,0,\pm1),\quad (-2,-4,\pm1),\\
 &(-2,8,\pm7),\quad (-2,-12,\pm7),\\
 &(-2,56,\pm41),\quad (-2,-60,\pm41),\\
 &(-2,336,\pm239),\quad (-2,-340,\pm239),\ldots.
\end{align*}

\section{Conclusion}

Theorem~\ref{thm:main} is an exact and duplication-controlled global
classification: choose \(m\mid y^2+2\), test whether
\(F_m(y)=my+(y^2+2)/m\) is a nonzero signed square, and the sign of that
square determines the sign of \(x\).  Theorem~\ref{thm:moduli} gives
the precise prime-factor obstruction on \(m=|x|\).  For fixed \(m\),
Proposition~\ref{prop:normform}, Theorem~\ref{thm:squarem}, and
Theorem~\ref{thm:pellorbits} provide an effective exhaustive procedure
in every case.  Finally, Theorem~\ref{thm:infinitefamily} proves
unconditionally that the full equation has infinitely many integer
solutions.

\end{document}
